% =====================================================================================
% Section 6 -- RQ2: calibration, its dependence on the question, and its uses (punchline 3 of
% docs/notes/narrative.md). Every number quoted here is in tables/calibration.tex (the paper's question) or
% tables/question-ablation.tex (S1Rank's question against ours).
% =====================================================================================
\section{RQ2: Are Decision Models Better Calibrated?}
\label{sec:calibration}

% Generated by `python -m cdm.paper_tables` (make tables, in paper/) from the rerank rows of
% results/results.csv -- do not edit by hand.
\begin{table}[!t]
\caption{Calibration on the judged pairs of TREC DL 2019--2023 (means; lower is better). ECE: expected
calibration error (15 bins); Brier: mean squared error of the probability against the 0/1 label. Per
foundation model: base model, decision fine-tunes in their better-calibrated mode (j/i: joint/isolated),
reranker fine-tunes; best in bold, second best underlined. Rows shaded as in Table~\ref{tab:triplets}.}
\label{tab:calibration}
\centering
\scriptsize
\setlength{\tabcolsep}{4pt}
\begin{tabular}[t]{lrr}
\toprule
System & ECE & Brier \\
\midrule
\multicolumn{3}{l}{\textit{Qwen3-0.6B}} \\
\rowcolor{rowbase}Qwen3-0.6B & 0.750 & 0.750 \\
\rowcolor{rowdecision}Kev-0.6B (i) & \underline{0.579} & \underline{0.543} \\
\rowcolor{rowreranker}Qwen3-Reranker-0.6B & \textbf{0.511} & \textbf{0.481} \\
\midrule
\multicolumn{3}{l}{\textit{Qwen3-8B}} \\
\rowcolor{rowbase}Qwen3-8B & 0.383 & 0.380 \\
\rowcolor{rowdecision}Kev-8B (j) & \textbf{0.243} & \underline{0.248} \\
\rowcolor{rowreranker}Qwen3-Reranker-8B & \underline{0.257} & \textbf{0.238} \\
\midrule
\multicolumn{3}{l}{\textit{Qwen3.5-0.8B}} \\
\rowcolor{rowbase}Qwen3.5-0.8B & 0.716 & 0.694 \\
\rowcolor{rowdecision}Kev-0.8B (i) & \underline{0.357} & \underline{0.309} \\
\rowcolor{rowreranker}Prism-0.8B & \textbf{0.243} & \textbf{0.193} \\
\midrule
\multicolumn{3}{l}{\textit{Qwen3.5-4B}} \\
\rowcolor{rowbase}Qwen3.5-4B & 0.296 & 0.269 \\
\rowcolor{rowdecision}Kev-4B (i) & \textbf{0.216} & \underline{0.207} \\
\rowcolor{rowreranker}Prism-4B & \underline{0.218} & \textbf{0.172} \\
\bottomrule
\end{tabular}
\hspace{2em}
\begin{tabular}[t]{lrr}
\toprule
System & ECE & Brier \\
\midrule
\multicolumn{3}{l}{\textit{Qwen3.5-9B}} \\
\rowcolor{rowbase}Qwen3.5-9B & 0.283 & 0.265 \\
\rowcolor{rowdecision}Kev-9B (j) & 0.237 & 0.209 \\
\rowcolor{rowdecision}Clef-Flash (j) & \underline{0.140} & \underline{0.165} \\
\rowcolor{rowreranker}Prism-9B & 0.242 & 0.183 \\
\rowcolor{rowreranker}\textit{Ours-9B} & \textbf{0.111} & \textbf{0.157} \\
\midrule
\multicolumn{3}{l}{\textit{Qwen3.8-27B}} \\
\rowcolor{rowbase}Qwen3.8-27B & 0.166 & 0.180 \\
\rowcolor{rowdecision}Kev-27B (j) & \underline{0.156} & \underline{0.159} \\
\rowcolor{rowdecision}Clef (j) & \textbf{0.147} & \textbf{0.157} \\
\midrule
\multicolumn{3}{l}{\textit{Other systems}} \\
\rowcolor{rowdecision}\jev{} (j) & \textbf{0.224} & \textbf{0.187} \\
\rowcolor{rowreranker}monoT5-3B & 0.312 & 0.295 \\
\rowcolor{rowreranker}bge-reranker-v2-m3 & \underline{0.231} & \underline{0.228} \\
\rowcolor{rowreranker}Set-Encoder-large & 0.495 & 0.426 \\
\bottomrule
\end{tabular}
\end{table}


Calibration is the main promise of decision models. Table~\ref{tab:calibration} reports it on the judged pairs of TREC
DL 2019--2023, together with the two uses we test. Table~\ref{tab:question-ablation} repeats every run of a decision
or untuned model with S1Rank's question.

\paragraph{Asked our question, decision models are well calibrated.}
The untuned models of 0.6B and 0.8B parameters say ``yes'' to almost everything (ECE 0.750 and 0.617); from 4B on,
they are reasonably calibrated already (0.117 to 0.231). Decision fine-tuning improves on the untuned model on every
backbone, for example from 0.231 to 0.113 on Qwen3-8B. From 4B on, the decision fine-tunes are also better calibrated
than the published reranker fine-tunes (Kev-8B 0.113 vs.\ Qwen3-Reranker-8B 0.257); on the two smallest backbones they
are not. The best calibrated systems in the study are \jev{} (0.075) and Kev-27B (0.076). Decision fine-tuning is not
the only way there: our 9B reranker fine-tune, trained with a plain yes/no loss, is as well calibrated as Clef-Flash
(0.111 vs.\ 0.112). Our 27B reranker fine-tune is less well calibrated than Kev-27B (0.110).

\paragraph{The advantage depends on the question.}
S1Rank's question asks whether the passage is relevant and defines relevance more loosely: the passage ``answers the
query or contains information that directly addresses it''. Both wordings are reasonable, and S1Rank's is the one
under which \jev{} was first shown to be a good reranker. The ranking hardly moves between the two: the mean nDCG@10
of the decision fine-tunes changes by $-2.5$ to $+2.6$ points, and with S1Rank's question, too, the best decision
fine-tune of every backbone up to 9B is 4 points or more behind the reranker fine-tune. The answer to RQ1 does not
depend on the wording. Calibration moves a lot. With S1Rank's question, almost every model is worse calibrated, often
by a factor of two or three (\jev{} joint: ECE 0.224 instead of 0.075). The decision fine-tunes are then at best as
well calibrated as the best reranker fine-tune of their backbone (Kev-8B 0.243 vs.\ Qwen3-Reranker-8B 0.257, Kev-27B
0.156 vs.\ Ours-27B 0.110). One sentence of the prompt decides whether decision models are better calibrated than
reranker fine-tunes.

The untuned models gain as much from our wording as the decision fine-tunes (Qwen3-8B: 0.383 with S1Rank's question,
0.231 with ours). A question that matches the definition of relevance in the labels makes the probabilities better
calibrated, with or without decision fine-tuning. The user of a decision model has to put that definition into words,
and a small change in the words changes the calibration severalfold. A reranker fine-tune has no such knob: its
notion of relevance is fixed by its training labels.

\paragraph{Calibration does not make up for a worse ranking.}
If every passage with $p \ge 0.5$ is kept, a reranker fine-tune keeps the better set on every backbone up to 9B
(e.g.\ $F_1$@.5 of 51.0 for Prism-4B vs.\ 42.6 for Kev-4B), although it is the worse calibrated model on most of
them. The threshold is not the problem: for the decision models, 0.5 is close to the best threshold of each
collection (regret of 0.9 to 3.1 points). A well-placed threshold on a worse ranking still yields the worse set. Where
the rankings are equally good, calibration does pay off: at 27B, the better calibrated Kev-27B keeps the better set
(57.2 vs.\ 48.3), while our fine-tune is underconfident and keeps only 7.9 passages per query.

\paragraph{Counting relevant passages.}
A ranking cannot tell how many passages are relevant; the sum of calibrated probabilities can. Asked our question,
most decision models from 4B on count better than the median guess (8.6 passages off) and better than the published
reranker fine-tunes of their backbone (Kev-8B 7.6 vs.\ Qwen3-Reranker-8B 12.9). Decision fine-tuning is not what makes
this work: the untuned Qwen3.8-27B counts most accurately of all (5.8), and Ours-9B counts better than Clef-Flash (7.3
vs.\ 8.1). With S1Rank's question, no decision model beats the median guess by more than a fraction of a passage
(best: Clef joint, 8.4).

\paragraph{Calibration after the fact.}
Rerankers can also be calibrated after training. Temperature scaling fitted on MS~MARCO dev lowers the ECE of the
worst calibrated rerankers (Qwen3-Reranker-0.6B from 0.511 to 0.326) and changes the others little. Platt scaling
backfires at the fixed cutoff ($F_1$@.5 of 8.1 or less): MS~MARCO dev has about one relevant passage per query, so the
fitted probabilities are far too low for TREC DL, where queries have many more. A calibration holds only for the
notion of relevance and the base rate it was fitted on (Sect.~\ref{sec:discussion}).

% After the text, so that LaTeX does not stack it with Table 3 on one page (that page overflows).
% Generated by `python -m cdm.paper_tables` (make tables, in paper/) from the rerank and `cdm cutoff` rows of
% results/results.csv -- do not edit by hand.
\begin{table*}[!t]
\caption{The question ablation: S1Rank's question (\emph{Rel.}: ``Is document \{doc\} relevant to the
query? Relevant means it answers the query or contains information that directly addresses it.'') against
the question used in the rest of the paper (\emph{Ans.}: ``Does document \{doc\} answer the query?''), for
every system that is asked a question. nDCG@10 ($\times 100$): mean over all ten collections (Avg) and over
TREC DL 2019--2023 (DL); ECE and Count as in Table~\ref{tab:calibration} (DL 2019--2023). $\Delta$: Ans.\ minus
Rel. Training: $\circ$~base model, $\blacktriangle$/$\triangle$~decision fine-tune (joint/isolated), $\blacksquare$~reranker fine-tune.}
\label{tab:question-ablation}
\centering
\scriptsize
\setlength{\tabcolsep}{3pt}
\begin{adjustbox}{max width=\textwidth}
\begin{tabular}{cllrrrrrrrrrrrr}
\toprule
 & & & \multicolumn{3}{c}{nDCG@10 Avg} & \multicolumn{3}{c}{nDCG@10 DL} & \multicolumn{3}{c}{ECE} & \multicolumn{3}{c}{Count} \\
\cmidrule(lr){4-6}\cmidrule(lr){7-9}\cmidrule(lr){10-12}\cmidrule(lr){13-15}
 & System & Mode & Rel. & Ans. & $\Delta$ & Rel. & Ans. & $\Delta$ & Rel. & Ans. & $\Delta$ & Rel. & Ans. & $\Delta$ \\
\midrule
\multicolumn{15}{l}{\textit{Qwen3-0.6B}} \\
\rowcolor{rowbase}$\circ$ & Qwen3-0.6B &  & 23.8 & 26.6 & $+$2.8 & 25.4 & 30.7 & $+$5.3 & 0.750 & 0.750 & \phantom{$+$}0.000 & 36.5 & 36.5 & \phantom{$+$}0.0 \\
\rowcolor{rowdecision}$\blacktriangle$ & Kev-0.6B & joint & 13.2 & 15.1 & $+$1.9 & 14.9 & 16.7 & $+$1.8 & 0.591 & 0.503 & $-$0.088 & 28.8 & 24.6 & $-$4.2 \\
\rowcolor{rowdecision}$\triangle$ & Kev-0.6B & isolated & 27.0 & 27.3 & $+$0.3 & 28.8 & 29.5 & $+$0.7 & 0.579 & 0.503 & $-$0.076 & 28.2 & 24.7 & $-$3.5 \\
\midrule
\multicolumn{15}{l}{\textit{Qwen3-8B}} \\
\rowcolor{rowbase}$\circ$ & Qwen3-8B &  & 53.3 & 51.1 & $-$2.2 & 56.9 & 54.0 & $-$2.9 & 0.383 & 0.231 & $-$0.152 & 17.5 & 8.7 & $-$8.8 \\
\rowcolor{rowdecision}$\blacktriangle$ & Kev-8B & joint & 38.8 & 38.7 & $-$0.1 & 50.9 & 49.3 & $-$1.6 & 0.243 & 0.113 & $-$0.130 & 12.8 & 7.6 & $-$5.2 \\
\rowcolor{rowdecision}$\triangle$ & Kev-8B & isolated & 52.7 & 51.7 & $-$1.0 & 54.2 & 53.8 & $-$0.4 & 0.283 & 0.154 & $-$0.129 & 14.1 & 8.7 & $-$5.4 \\
\midrule
\multicolumn{15}{l}{\textit{Qwen3.5-0.8B}} \\
\rowcolor{rowbase}$\circ$ & Qwen3.5-0.8B &  & 37.6 & 33.9 & $-$3.7 & 40.3 & 39.1 & $-$1.2 & 0.716 & 0.617 & $-$0.099 & 34.8 & 29.9 & $-$4.9 \\
\rowcolor{rowdecision}$\blacktriangle$ & Kev-0.8B & joint & 25.3 & 27.9 & $+$2.6 & 28.9 & 30.8 & $+$1.9 & 0.395 & 0.398 & $+$0.003 & 19.5 & 19.6 & $+$0.1 \\
\rowcolor{rowdecision}$\triangle$ & Kev-0.8B & isolated & 36.0 & 34.2 & $-$1.8 & 36.1 & 35.3 & $-$0.8 & 0.357 & 0.336 & $-$0.021 & 18.0 & 16.9 & $-$1.1 \\
\midrule
\multicolumn{15}{l}{\textit{Qwen3.5-4B}} \\
\rowcolor{rowbase}$\circ$ & Qwen3.5-4B &  & 52.7 & 52.4 & $-$0.3 & 55.8 & 56.1 & $+$0.3 & 0.296 & 0.181 & $-$0.115 & 14.7 & 9.8 & $-$4.9 \\
\rowcolor{rowdecision}$\blacktriangle$ & Kev-4B & joint & 49.8 & 48.2 & $-$1.6 & 53.4 & 51.9 & $-$1.5 & 0.251 & 0.201 & $-$0.050 & 13.2 & 11.6 & $-$1.6 \\
\rowcolor{rowdecision}$\triangle$ & Kev-4B & isolated & 51.3 & 49.9 & $-$1.4 & 52.7 & 51.8 & $-$0.9 & 0.216 & 0.132 & $-$0.084 & 11.5 & 8.4 & $-$3.1 \\
\midrule
\multicolumn{15}{l}{\textit{Qwen3.5-9B}} \\
\rowcolor{rowbase}$\circ$ & Qwen3.5-9B &  & 53.7 & 53.1 & $-$0.6 & 56.7 & 56.6 & $-$0.1 & 0.283 & 0.165 & $-$0.118 & 14.1 & 8.7 & $-$5.4 \\
\rowcolor{rowdecision}$\blacktriangle$ & Kev-9B & joint & 50.2 & 48.3 & $-$1.9 & 54.0 & 51.9 & $-$2.1 & 0.237 & 0.160 & $-$0.077 & 12.2 & 9.5 & $-$2.7 \\
\rowcolor{rowdecision}$\triangle$ & Kev-9B & isolated & 50.1 & 49.9 & $-$0.2 & 52.8 & 52.2 & $-$0.6 & 0.348 & 0.221 & $-$0.127 & 16.9 & 11.8 & $-$5.1 \\
\rowcolor{rowdecision}$\blacktriangle$ & Clef-Flash & joint & 52.8 & 51.8 & $-$1.0 & 57.0 & 56.6 & $-$0.4 & 0.140 & 0.112 & $-$0.028 & 8.9 & 8.1 & $-$0.8 \\
\rowcolor{rowdecision}$\triangle$ & Clef-Flash & isolated & 54.7 & 53.8 & $-$0.9 & 58.9 & 57.8 & $-$1.1 & 0.173 & 0.144 & $-$0.029 & 9.7 & 8.2 & $-$1.5 \\
\midrule
\multicolumn{15}{l}{\textit{Qwen3.8-27B}} \\
\rowcolor{rowbase}$\circ$ & Qwen3.8-27B &  & 57.4 & 56.4 & $-$1.0 & 59.9 & 59.6 & $-$0.3 & 0.166 & 0.117 & $-$0.049 & 8.9 & 5.8 & $-$3.1 \\
\rowcolor{rowdecision}$\blacktriangle$ & Kev-27B & joint & 61.1 & 60.1 & $-$1.0 & 63.5 & 63.3 & $-$0.2 & 0.156 & 0.076 & $-$0.080 & 8.5 & 6.0 & $-$2.5 \\
\rowcolor{rowdecision}$\triangle$ & Kev-27B & isolated & 59.8 & 59.0 & $-$0.8 & 62.5 & 62.3 & $-$0.2 & 0.262 & 0.180 & $-$0.082 & 12.9 & 9.4 & $-$3.5 \\
\rowcolor{rowdecision}$\blacktriangle$ & Clef & joint & 56.1 & 54.8 & $-$1.3 & 61.1 & 60.6 & $-$0.5 & 0.147 & 0.091 & $-$0.056 & 8.4 & 6.7 & $-$1.7 \\
\rowcolor{rowdecision}$\triangle$ & Clef & isolated & 56.4 & 53.9 & $-$2.5 & 58.5 & 58.3 & $-$0.2 & 0.240 & 0.153 & $-$0.087 & 12.1 & 8.3 & $-$3.8 \\
\midrule
\multicolumn{15}{l}{\textit{Other systems}} \\
\rowcolor{rowdecision}$\blacktriangle$ & \jev{} & joint & 60.5 & 59.1 & $-$1.4 & 62.7 & 62.2 & $-$0.5 & 0.224 & 0.075 & $-$0.149 & 11.3 & 6.3 & $-$5.0 \\
\rowcolor{rowdecision}$\triangle$ & \jev{} & isolated & 59.5 & 57.1 & $-$2.4 & 61.4 & 60.0 & $-$1.4 & 0.260 & 0.125 & $-$0.135 & 12.9 & 7.8 & $-$5.1 \\
\rowcolor{rowdecision}$\triangle$ & \laya{} & isolated & 44.0 & 43.5 & $-$0.5 & 47.3 & 47.0 & $-$0.3 & 0.131 & 0.127 & $-$0.004 & 9.3 & 9.2 & $-$0.1 \\
\bottomrule
\end{tabular}
\end{adjustbox}
\end{table*}

